\section{Tool-Vergleich und Fazit}
\label{sec:fazit}

\subsection{Tool-Vergleich: Celonis vs.\ Python (PM4Py)}
\label{sec:toolvergleich}

Die parallele Untersuchung mit der kommerziellen Plattform Celonis und der Open-Source-Bibliothek PM4Py zeigt, dass sich beide Werkzeuge komplementär ergänzen, in der Datenverarbeitung und bei einzelnen Kennzahlen jedoch methodisch bedingte Abweichungen aufweisen (\autoref{tab:toolvergleich}).

\begin{table}[htbp]
\centering
\caption{Systematischer Vergleich von Celonis und Python (PM4Py) im Analyseprozess}
\label{tab:toolvergleich}
\footnotesize
\renewcommand{\arraystretch}{1.15}
\begin{tabularx}{\textwidth}{@{}p{2.9cm} X X@{}}
\toprule
\textbf{Dimension} & \textbf{Celonis (Process Analytics \& PQL)} & \textbf{Python (PM4Py)} \\
\midrule
\textbf{Datenimport \& Attribute} &
Direkter XES-Import erzeugt nur eine Aktivitätstabelle; Trace-Attribute (\texttt{Amount}, \texttt{Permit ID}) fehlen ohne vorherige CSV-Konvertierung. &
Voller Zugriff auf alle Trace- und Event-Attribute; deckt 448 Dummy-Permits (\texttt{travel permit 423}) und 3.078 redundante Event-IDs auf. \\
\textbf{Exploration \& Vorverarbeitung} &
Sehr schnelle, interaktive Filterung über Varianten, Zeiträume und Einzelfälle (\textit{Case Explorer}); PQL erlaubt flexible KPI-Dashboards. &
Hohe algorithmische Freiheit: Text-Clustering der 37 Eventnamen (Exact Cover / Algorithmus X), Entfernen der Trip-Events und Zeitsplits. \\
\textbf{Process Discovery \& Conformance} &
Frequenzbasierter Process Explorer und BPMN-Abgleich (Variante 1 vs.\ Top-3 mit XOR); visualisiert den Strukturbruch 2018 unmittelbar. &
Systematisches Benchmarking vieler Miner (Heuristic, Inductive, ILP); exakte Fitness-, Precision- und F1-Berechnung (Anstieg von 0,902 auf 0,961). \\
\textbf{Ergebnisvergleich (Kennzahlen)} &
Varianten: 100 (Dom.) / 756 (Int.).\newline
Ø Durchlaufzeit Int.: weist je nach View 66\,d (Happy Path), 75\,d (getrimmt) oder 86\,d (alle Cases) aus. &
Varianten: 99 (Dom.) / 753 (Int.) (Tie-Breaking bei gleichen Zeitstempeln).\newline
Ø Durchlaufzeit Int.: 86\,d auf Rohdaten. \\
\bottomrule
\end{tabularx}
\end{table}

\subsection{Fazit und Limitationen}
\label{sec:zusammenfassung}
Trotz prozessualer Veränderungen im Zeitverlauf, konnten bereits mit den Rohdaten in Celonis und in Python mögliche Prozessabläufe identifiziert werden, die zu einer starken Abdeckung der Gesamtlogs führen.
Mit F1-Scores von über 0,9 konnte konsequent bestätigt werden, dass die Prozessmodelle die beobachtbaren Event-Logs abbildeten.
Mithilfe von Textanalyse und Exact-Cover-Clustering konnten zudem mögliche Objektzentrierungen und Aktivitäten-Clustering stattfinden, was die Interpretation stütze.\par
Es lässt sich festhalten, dass für inländische (Domestic) Reisen eine Reisekostenabrechnung (Declaration) vorangeht, um die Bezahlung nach einem mehrstufigen Genehmigungsprozess freizugeben.
Für internationale (International) Reisen wird hingegen auch ein Reiseantragsgenehmigungsprozess (Permit) vorausgesetzt, um die Reisekostenabrechnung loszutreten. Zusätzlich werden im International-Log auch Reminder verschickt im Gesamtprozess und auch die Reisedaten (Trip) erfasst.
Auffällig ist, dass in dem internationalen Log identische Events mehrfach vorkommen und "wiederverwendet" werden.
Auch gibt es Fälle, in denen in den Trace-Metadaten -- welche belastbar ausführlicher sind als beim Domestic-Log -- Permit-Metadaten anders befüllt werden und in dieser Abweichung ein Zusammenhang mit International-Traces ohne Reisegenehmigung (Permit) vorliegt.
In diesen Fällen konnte eine Reisekostenabrechnung (Declaration) ohne Reisegenehmigung (Permit) stattfinden, was auf eine mögliche Prozessabweichung hindeutet und zunächst nicht als konform wirkt in Rücksicht auf den geringen Trace-Anteil.
\par
Cash-Flows (Payment) bedingen im Allgemeinen den Abschluss des Genehmigungsprozesses komplett. In beiden Logs finden sich dafür aber auch Gegenbeispiele, bei denen das Payment durchgeführt wurde, bevor eine volle Genehmigung vorlag.
Im Log waren zudem auch nicht nur konforme positiv abgeschlossene Fälle enthalten, sondern auch 1301 (Domestic) bzw. 1755 (International) abgelehnte Fälle, bei denen im Log keine weitere Bearbeitung erfolgte.
Bei diesen Traces ist von einer finalen Ablehnung auszugehen. Insgesamt konnte der Grundprozess annähernd präzise abgebildet werden, wobei festgestellt wird, dass weitere Metainformationen zur Prozesshistorie und Organisationen unterstützen würden.
